\section{Scope, contributions, and non-claims}
\label{sec:scope}

This paper develops an \emph{alternative deterministic substrate theory} whose goal is to present a consistent
continuum formulation of its main building blocks and to make its open problems explicit.
The focus is the continuous model; discrete simulations are not treated as fundamental and appear only as an
implementation appendix.

\subsection{Core building blocks}
The theory fragment developed here is organized around four coupled elements:
\begin{enumerate}
\item \textbf{Continuum 3D brane in 4D embedding.}
A rotationally invariant hyperelastic action on an intrinsic $3$-manifold $\Omega\subset\R^3$ determines the
dynamics of an embedding $\vecX:\Omega\times\R\to\R^4$ (\Cref{sec:continuum-model}).

\item \textbf{Ordered narrowband carrier sector.}
We assume the substrate admits an ordered, narrowband carrier sector (taken to be primarily longitudinal)
that provides phase rigidity and enables a controlled multiscale envelope description
(\Cref{sec:geophase,sec:effective-field}).

\item \textbf{Berry/Wilczek--Zee connections as gauge structure.}
Adiabatic transport of a polarization subspace induces a Berry connection (rank-1) or a Wilczek--Zee
connection (non-Abelian).
These connections are interpreted as the natural origin of gauge degrees of freedom in a coarse-grained
description, and provide the bridge to a $U(1)$-like gauge picture and its non-Abelian extension
(\Cref{sec:geophase,sec:effective-field}).

\item \textbf{Solitonic particles and a Dirac-envelope split.}
Particle-like states are localized solitons whose spatial structure is organized by spherical symmetry and
spherical-harmonic mode families (\Cref{sec:localized}).
Conceptually, we separate (i) particle-internal structure (topology/geometry and its internal labels, including
matter/antimatter distinction and spin state) from (ii) the underlying local phase structure of the carrier sector.
The standard Dirac equation mixes both; here it is treated as an effective envelope description that couples the
particle sector to the local phase/gauge structure.
The spin-$\tfrac12$ $4\pi$ aspect is attributed to geometric phase/holonomy of an internal polarization subspace.
\end{enumerate}

\subsection{Two working hypotheses to be discriminated}
To keep the proposal testable, we distinguish two nonexclusive working hypotheses that can be separated by
numerical experiments:
\begin{description}
\item[Hypothesis~A (transverse-channel confinement).]
Localization is dominated by nonlinear transverse--lateral coupling of the embedded brane
(\Cref{subsec:gravity-channel,sec:localized}); gauge variables extracted from Berry/Wilczek--Zee holonomy are
primarily \emph{readouts} of the ordered carrier sector and need not correlate strongly with confinement strength.

\item[Hypothesis~B (holonomy-tied confinement).]
Localization relies on a near-degenerate internal polarization subspace whose transport induces a robust Berry/
Wilczek--Zee connection; in this case confinement and the extracted holonomy/curvature diagnostics should correlate
systematically, and a non-Abelian ($n>1$) description may be required.
\end{description}

\subsection{Interpretive stance on quantum phenomena}
We do not treat the probabilistic postulates of quantum theory as fundamental.
Apparent point-like interactions and discrete outcomes are attributed to deterministic wave dynamics in a
narrowband medium, together with Fourier-localization constraints and threshold-like nonlinear exchange events.
This stance motivates how localized solitons can exhibit particle-like scattering while remaining wave states.

\subsection{Non-claims (what is \emph{not} asserted here)}
\begin{enumerate}
\item We do not claim a complete derivation of the Standard Model, nor a finished replacement for QED/QFT.
\item We do not claim that the geometric connections are \emph{identical} to electromagnetism; rather, we present a
consistent identification of gauge degrees of freedom with Berry/Wilczek--Zee holonomy in a narrowband sector and
develop its continuum consequences.
\item We do not claim a unique constitutive law; we specify a concrete isotropic hyperelastic choice
(St.\ Venant--Kirchhoff) as a baseline while keeping the constitutive layer modular.
\item We do not claim the discrete lattice approximation is fundamental; it is only a numerical realization of the
continuum theory and kept out of the main narrative.
\end{enumerate}
